\clearpage\thispagestyle{empty}
\pdfbookmark[0]{Quellen zum Kometenbericht vom Dezember 1807}{comet1807december-sources}
\begin{center}\textbf{Quellen zum Kometenbericht vom Dezember 1807}\end{center}
\noindent Diese KI-Übertragung in heutiges Deutsch folgt dem Kometenauszug vom Dezember 1807 in \href{https://gdz.sub.uni-goettingen.de/id/PPN236007467}{Gauß’ \textit{Werken}, Band~VI} (Göttingen, 1874), S.~292. Franz Xaver von Zach leitet den Brief von Carl Friedrich Gauß vom 16. November 1807 ein. Der Erstdruck erschien in Beitrag XLIV, \textit{Fortgesetzte Nachrichten über den Cometen von 1807}, in der \textit{Monatlichen Correspondenz zur Beförderung der Erd- und Himmelskunde}, Band~XVI, Dezember 1807, \href{https://www.digitale-sammlungen.de/en/view/bsb10538607?page=576}{S.~562}--\href{https://www.digitale-sammlungen.de/en/view/bsb10538607?page=581}{567}. Die Einleitung zum Brief beginnt unten auf \href{https://www.digitale-sammlungen.de/en/view/bsb10538607?page=578}{S.~564}. Gauß’ Bahnelemente und seine abschließenden Bemerkungen stehen auf \href{https://www.digitale-sammlungen.de/en/view/bsb10538607?page=579}{S.~565}.
\par\medskip
Die Tabelle der Werkausgabe enthält die eingeklammerte Bahnneigung und schreibt in der ersten Zeile \textit{duchs}, während im Erstdruck \textit{durchs} steht. Die Zeitschriftendigitalisate stammen von der Bayerischen Staatsbibliothek. Übertragung und Satz wurden mit KI erstellt.
\clearpage
